\documentclass[11pt,a4paper]{article}
\usepackage[T1]{fontenc}
\usepackage[utf8]{inputenc}
\usepackage{lmodern}
\usepackage[margin=27mm]{geometry}
\usepackage{microtype}
\usepackage{amsmath,amssymb,amsthm,mathtools,mathrsfs}
\usepackage{booktabs,array,longtable}
\usepackage{enumitem}
\usepackage{xcolor}
\usepackage{graphicx}
\usepackage{tikz}
\usetikzlibrary{arrows.meta,positioning}
\usepackage{algorithm}
\usepackage{algpseudocode}
\usepackage[round,authoryear]{natbib}
\usepackage{aliascnt}
\usepackage[colorlinks=true,linkcolor=blue!50!black,citecolor=blue!50!black,urlcolor=blue!60!black]{hyperref}
\usepackage[capitalise,noabbrev]{cleveref}
\usepackage{fancyhdr}
% Shared mathematical notation. Arc i -> j means: i requires j.
% Siamo nel mondo dei grafi e delle closure: V sono i vertici, non "item".
\newcommand{\Vset}{\mathcal V}
\newcommand{\Aset}{\mathcal A}
\newcommand{\Gset}{\mathcal G}
\newcommand{\Cset}{\mathscr C}
\newcommand{\Rset}{\mathbb R}
\newcommand{\Sseq}{\mathcal S}
\newcommand{\Fset}{\mathcal F}
\newcommand{\Qset}{\mathcal Q}
\newcommand{\Bset}{\mathcal B}
\newcommand{\Nset}{\mathcal N}
\newcommand{\Tset}{\mathcal T}
\newcommand{\Kset}{\mathcal K}   % blocchi canonici K_1, ..., K_k
% Matrici e vettori in grassetto; le lettere greche via \boldsymbol.
\newcommand{\mat}[1]{\mathbf{#1}}
\newcommand{\vect}[1]{\mathbf{#1}}
\newcommand{\vecg}[1]{\boldsymbol{#1}}
\newcommand{\ind}{\chi}
\newcommand{\divg}{\operatorname{div}}
\newcommand{\supp}{\operatorname{supp}}
% Nome del metodo (decisione D18): minuscolo, sul modello di *network simplex*.
% \FS in mezzo alla frase, \FSC a inizio frase, \CS per la forma breve.
\newcommand{\FS}{ratio-closure simplex}
\newcommand{\FSC}{Ratio-closure simplex}
\newcommand{\CS}{closure simplex}
\newcommand{\LP}{\textsc{LP}}
\newcommand{\rhoT}{\rho}

\numberwithin{equation}{section}
\newtheorem{theorem}{Theorem}[section]
\newaliascnt{proposition}{theorem}
\newtheorem{proposition}[proposition]{Proposition}
\aliascntresetthe{proposition}
\newaliascnt{lemma}{theorem}
\newtheorem{lemma}[lemma]{Lemma}
\aliascntresetthe{lemma}
\newaliascnt{corollary}{theorem}
\newtheorem{corollary}[corollary]{Corollary}
\aliascntresetthe{corollary}
\crefname{proposition}{Proposition}{Propositions}
\crefname{lemma}{Lemma}{Lemmas}
\crefname{corollary}{Corollary}{Corollaries}
\theoremstyle{definition}
\newtheorem{definition}[theorem]{Definition}
\newtheorem{example}[theorem]{Example}
\theoremstyle{remark}
\newtheorem{remark}[theorem]{Remark}
\setlist{topsep=4pt,itemsep=2pt,parsep=0pt}
\setlength{\emergencystretch}{2em}
\pagestyle{fancy}
\fancyhf{}
\fancyhead[L]{\small A ratio-closure simplex on general precedence graphs}
\fancyhead[R]{\small Review draft}
\fancyfoot[C]{\thepage}
\setlength{\headheight}{14pt}
\title{A Ratio-Closure Simplex\\ for Maximum-Ratio Closure and its Canonical Decomposition\\[5pt]
\large A didactic derivation on general directed acyclic graphs}
\author{Working document for mathematical review}
\date{5 September 2026}
\hypersetup{
  pdftitle={A Ratio-Closure Simplex for Maximum-Ratio Closure on Precedence Graphs},
  pdfauthor={Working document for mathematical review},
  pdfsubject={A didactic derivation on general directed acyclic graphs},
  pdfkeywords={ratio-closure simplex, maximum-ratio closure, precedence graph, linear programming}
}

\begin{document}
\maketitle
\begin{abstract}
This tutorial develops the \FS{} for the maximum-ratio closure
subproblem underlying a linear program with arbitrary precedence constraints,
positive weights, profits of either sign, unit upper bounds, and one capacity
constraint. The central object is a spanning forest of \emph{nonbasic slack
variables}, together with roots for its zero components. We derive the basis
representation, oriented subtree formulas for reduced costs and directions,
the ratio test, all four structural pivot updates, and an exact-arithmetic
termination argument. This representation supports specialized computation
of both entering-variable prices and the pivot column: all reduced costs
take $O(n)$ arithmetic operations after the forest is available, and a full
pivot takes $O(n+m)$ with a complete arc scan. Repeated ratio optimization recovers the block
decomposition and an optimum of the capacitated LP. A complete eight-pivot
trace includes six degenerate pivots and is checked independently with rational
linear algebra. This is a mathematical and didactic draft, not a performance
report or a claim that the optimized C++ solver has already been implemented.
\end{abstract}

\paragraph{The three structural ingredients.}
A basis is encoded by a forest and its roots, with exactly one rootless
tree. This combinatorial object gives both the reduced costs used to choose
an entering variable and the corresponding pivot column through subtree
sums and signed component movements. These are specialized solutions of
the two basis systems of the revised simplex method. Sections~\ref{sec:forest} and
\ref{sec:directions} establish these three ingredients; Section~\ref{sec:implementation}
summarizes the arithmetic cost of each operation. The input precedence graph
may be any DAG: the forest represents the current basis.

\paragraph{Status and reading guide.}
The notation and general LP follow the supplied manuscript
\citet{DoseFuriniLocatelli2026}. The forest viewpoint starts from the internal
notes \citet{ForestNotes}; here its formulas are derived for arbitrary arc
orientations. Sections~\ref{sec:model}--\ref{sec:ratio} introduce only the
background needed by the algorithm. Sections~\ref{sec:forest}--\ref{sec:trace}
are the core: basis, pricing, directions, pivots, and a worked execution.
Section~\ref{sec:outer} returns to the original LP.
The proofs concern exact arithmetic. The numerical and performance variants
in Section~\ref{sec:implementation} are implementation proposals, not measured
results. The checks in Section~\ref{sec:verification} support review of the
formulas; they do not replace the proofs or constitute a production solver.
No novelty claim is made before comparison with the relevant simplex and
closure literature.

\tableofcontents
\clearpage
\section{The linear program and the running instance}\label{sec:model}

Let $\Gset=(\Vset,\Aset)$ be a directed acyclic graph (DAG), with
$n=|\Vset|$ vertices and $m=|\Aset|$ arcs. An arc $(i,j)$ means that
\emph{vertex $j$ is a prerequisite of vertex $i$}. Its inequality is therefore
$x_i\le x_j$. This convention is used in every formula and diagram below.
Each vertex has profit $p_i\in\Rset$ and weight $w_i>0$. There is a single vertex
type: no bipartition or alternation of vertex types is assumed, a vertex may have
prerequisites and itself be a prerequisite, and profits may be negative, so
selecting a profitable vertex can require selecting loss-making predecessors.

\begin{definition}[Closure]
A set $C\subseteq\Vset$ is a closure if $i\in C$ and $(i,j)\in\Aset$
imply $j\in C$. Write $\Cset$ for all closures, including the empty set.
For $S\subseteq\Vset$, let
\[
 p(S)=\sum_{i\in S}p_i,\qquad w(S)=\sum_{i\in S}w_i,
 \qquad \varrho(S)=\frac{p(S)}{w(S)}\quad(S\ne\emptyset).
\]
\end{definition}

\subsection{Names and notation, taken from the literature}\label{sec:names}
Every term used below already has an owner, and this tutorial uses the owner's
word rather than a new one.

\begin{center}
\begin{tabular}{@{}lll@{}}
\toprule
Object & Name used here & Source \\
\midrule
$\Gset=(\Vset,\Aset)$ & digraph, vertices, arcs & standard \\
$C\subseteq\Vset$ down-closed & \emph{closure} & \citet{Picard1976} \\
$\max_C p(C)/w(C)$ & \emph{maximum-ratio closure} & \citet{Hochbaum2024} \\
the same set, order-theoretic & $\rho$-maximal \emph{initial set} & \citet{Sidney1975,Lawler1978} \\
the nested chain of closures & \emph{Sidney decomposition} & \citet{Sidney1975} \\
its successive differences $\Kset_r$ & \emph{blocks} & this tutorial \\
forest with one distinguished vertex per tree & \emph{rooted} forest, \emph{roots} & \citet{Yang2026} \\
the tree with no such vertex & \emph{rootless} tree & \citet{Yang2026} \\
forest with per-branch aggregates & \emph{normalized tree} & \citet{LerchsGrossmann1965} \\
\bottomrule
\end{tabular}
\end{center}

\noindent
\citet{Picard1976} writes \emph{maximal} closure where current usage writes
\emph{maximum} closure; the latter is used here. Matrices and vectors are set in
bold, scalars and indexed entries are not, so $\vect{w}$ is the weight vector
and $w_i$ one of its entries.

\subsection{The problem: maximum-ratio closure}\label{sec:problem}
The \FS{} solves one problem, and it is a \emph{ratio} problem over closures:
\begin{equation}\label{eq:ratioproblem}
  \rho^* \;=\; \max\Bigl\{\,\varrho(C)=\frac{p(C)}{w(C)}
              \;:\; \emptyset\ne C\subseteq\Vset,\ C\in\Cset \,\Bigr\}.
\end{equation}
Note what is absent: there is no capacity, and no upper bound on any variable.
The name of the method follows \eqref{eq:ratioproblem}, in the way that
\emph{network simplex} names the class of flow problems rather than the spanning
tree it maintains.

Two devices linearize \eqref{eq:ratioproblem}. Normalizing the denominator
\citep{CharnesCooper1962} gives a single linear program, stated as
\eqref{eq:normalized} and proved equivalent in \cref{thm:normalization}; that
program, and its structure, is what the rest of this tutorial works on, and
\cref{sec:class} identifies the class it belongs to. Treating the ratio
parametrically instead \citep{Dinkelbach1967} turns each iterate into a maximum
closure and hence a minimum cut \citep{Picard1976}; that is the route of the
comparison methods in \cref{app:flow}, not of the method developed here.

\subsection{Application: adding a capacity row}\label{sec:capacity}
The program \eqref{eq:ratioproblem} is an oracle, and the problem it serves is the
LP relaxation of a precedence-constrained knapsack. For capacity $c\ge0$,
\begin{subequations}\label{eq:lp}
\begin{align}
 z(c)=\max\quad &\sum_{i\in\Vset}p_i x_i,\label{eq:lpobj}\\
 \text{s.t.}\quad &\sum_{i\in\Vset}w_i x_i\le c,\label{eq:lpcap}\\
 &x_i-x_j\le0 &&(i,j)\in\Aset,\label{eq:lpprec}\\
 &0\le x_i\le1 &&i\in\Vset.\label{eq:lpbox}
\end{align}
\end{subequations}
The indicator $\ind^C$ is feasible for the precedence and box constraints
exactly when $C$ is a closure; the capacity row makes the LP generally
fractional. Solving \eqref{eq:lp} for a given $c$ needs no new machinery beyond
\eqref{eq:ratioproblem}: repeated calls of the oracle on residual graphs produce
closures of decreasing ratio, and $c$ decides how far along that sequence one
walks and which fraction of the critical block is taken.
\Cref{sec:outer} does this, and \cref{thm:outer} gives the whole function
$z(\cdot)$ at once, piecewise linear with one breakpoint per block, at no
extra cost. The capacity therefore never appears inside the simplex; it appears
only in the procedure that calls it, which is why the method is named after
\eqref{eq:ratioproblem} and not after the knapsack.

The same machinery answers a question in which no capacity appears at all.
Driving the oracle to exhaustion --- extract a block, delete it, ask again on
the residual until nothing is left --- produces the \emph{canonical
decomposition}: the complete sequence $\Kset_1,\dots,\Kset_k$ of blocks of
strictly decreasing ratio, together with the threshold of every vertex and the
breakpoints of the parametric closure value. This is the object the method is
built to deliver, and it is what the solver computes under the tasks
\texttt{full-path} and \texttt{positive-path}; a single capacity is then just a
query against a sequence already in hand. The sequence is not a device of this
tutorial: it is the decomposition of \citet{Sidney1975}, and the same object is
studied on precedence forests, under the name \emph{closure layers}, by the
companion work of this project. What is specific here is that a simplex basis,
rather than a flow, is what walks it.

\subsection{Eight vertices, several precedence levels}
We use the running instance from the supplied manuscript. Its data are
stored once in \texttt{tests/fixtures/paper\_example8.json}; the tables and
figures below are generated from that file.
\begin{table}[htbp]
\centering\input{../build/tutorial/generated/instance_table.tex}
\caption{Profits and weights of the running instance.}\label{tab:data}
\end{table}
Its arcs, in the fixed order used to number the slack variables, are
\begin{equation}\label{eq:arcs}
\begin{split}
\Aset=\{&(5,1),(5,2),(5,6),(6,3),(7,3),\\
         &(7,4),(8,5),(8,6),(8,7)\}.
\end{split}
\end{equation}
For example, vertex 8 requires vertices 5, 6 and 7, which in turn impose other
prerequisites. At $c=4$, the closure $\{1,3,6\}$ is an optimal integer
selection of value 5. The LP value will be 7. The integer problem is used
only to illustrate the difference; it is not the problem solved here.
\begin{figure}[htbp]
\centering\input{../build/tutorial/generated/graph.tex}
\caption{The general precedence graph. A label $(p_i,w_i)$ is attached to
each vertex; arrows point to prerequisites.}\label{fig:graph}
\end{figure}

\subsection{The dual and a convention to remember}
For arc multipliers $\alpha$, use the divergence convention of
\citet{DoseFuriniLocatelli2026}:
\[
\divg_i(\alpha)=\sum_{j:(i,j)\in\Aset}\alpha_{ij}
                  -\sum_{j:(j,i)\in\Aset}\alpha_{ji}.
\]
The dual of \eqref{eq:lp} is
\begin{equation}\label{eq:lpdual}
\begin{aligned}
\min\quad &c\lambda+\sum_i\mu_i\\
\text{s.t.}\quad &w_i\lambda+\mu_i+\divg_i(\alpha)\ge p_i &&i\in\Vset,\\
&\lambda\ge0,\quad\mu_i\ge0,\quad\alpha_{ij}\ge0.
\end{aligned}
\end{equation}
The multiplier $\lambda$ prices capacity; $\mu_i$ prices the upper bound.
Later, $s_{ij}$ will denote a \emph{primal slack}, not an arc multiplier.
Keeping $s_{ij}$, $\alpha_{ij}$ and $\lambda$ distinct prevents a collision
with the notation in the historical forest notes.

\subsection{Where this sits in the literature}\label{sec:prior}
Nothing in \eqref{eq:ratioproblem} or in the forest structure that follows from it is
new, and it is better to say so here than to let a reader discover it later.
The companion document in \path{literature/} does this carefully; the four
blocks are:

\begin{description}[leftmargin=*,style=nextline]
\item[Closure and its reduction to flow.] Maximum closure and its solution by
  minimum cut are classical \citep{Picard1976,PicardQueyranne1982}; the
  pseudoflow algorithm \citep{Hochbaum2008} is what a flow-based competitor
  would use today, and \citet{HochbaumQueyranne2003} place the closure problem
  and isotonic regression in one family.
\item[The ratio objective and the decomposition.] The normalization is
  \citet{CharnesCooper1962}; the parametric reading is
  \citet{Dinkelbach1967}; the decomposition into blocks of decreasing ratio is
  \citet{Sidney1975} and \citet{Lawler1978}; parametric flow over the whole
  ratio range is \citet{GalloEtAl1989} and \citet{Eppstein2018}, and the
  divide-and-conquer scheme is \citet{EisnerSeverance1976}. The union
  convention that makes a block canonical is the maximal source side of a
  minimum cut \citep{PicardQueyranne1980}.
\item[Forests carrying branch aggregates.] This is the oldest and closest
  antecedent. The open-pit design algorithm of \citet{LerchsGrossmann1965}
  maintains a \emph{normalized tree}, carries the mass of each branch,
  classifies an arc by the sign of the mass it supports, and pivots by adding
  one arc and removing a branch's root arc; \citet{Hochbaum2001} sets this out
  and recovers a flow from the tree in linear time, and also gives a parametric
  variant with the cost of a single run. Substituting $b_i=p_i-\rho w_i$ turns
  the strong/weak test into the sign of the reduced costs derived in
  \cref{sec:directions}.
\item[Recent simplex work on the same shapes.] \citet{Yang2026} develops the
  revised-simplex algebra on rooted spanning forests for
  budget-constrained total-variation problems, which by reading (c) of
  \cref{sec:class} is the same structural family as \eqref{eq:normalized}.
  \citet{Wuille2025} states \eqref{eq:normalized} itself, with $n-1$ nonbasic
  variables, a spanning forest of
  tight dependency slacks, at most one free vertex variable per component and a
  single positive component; \citet{MollakhaniEtAl2026} formalize that as a
  spanning-forest linearization algorithm. Simplex methods on an incidence
  matrix with one side row are classical
  \citep{AhujaEtAl1993,Schrijver2003,HolzhauserEtAl2017}.
\end{description}

What this tutorial develops is therefore a \emph{specialization}: the ordinary
primal simplex carried through exactly on \eqref{eq:normalized}, with the closed
forms that specialization produces. It is not a new class of basis and not a
new paradigm.

\subsection{Cycles and limits of the model}
If the input digraph contains a directed cycle, the corresponding inequalities
force equality of the variables on that cycle. Contract each strongly
connected component, sum its profits and weights, and retain a map to expand
the final solution. The resulting graph is a DAG with positive aggregated
weights. Thus the DAG assumption is a preprocessing convention for general
directed precedence graphs. In contrast, zero or negative weights, several
independent capacity rows, heterogeneous upper bounds, and inequalities of
the form $x_i\le\gamma_{ij}x_j$ require additional analysis. They are not
included in the statements of this draft.

\section{The ratio oracle and its normalization}\label{sec:ratio}

Suppose a closed prefix $M$ has already been selected. Put
$R=\Vset\setminus M$ and retain only arcs whose two endpoints lie in $R$.
A set $S\subseteq R$ is a closure in this residual graph if and only if
$M\cup S$ is a closure in the original graph. Arcs from $R$ to $M$ already
have their prerequisites satisfied. The next subproblem is
\begin{equation}\label{eq:ratio}
 \rho^*(R)=\max\left\{\frac{p(S)}{w(S)}:
              \emptyset\ne S\subseteq R,\ S\text{ a residual closure}\right\}.
\end{equation}
We now work on one nonempty residual graph, suppress the superscript $R$,
and use $n,m$ for its size.

\subsection{Normalization removes the quotient}
The \FS{} solves the following LP:
\begin{equation}\label{eq:normalized}
\begin{aligned}
\max\quad &\sum_i p_i y_i\\
\text{s.t.}\quad &\sum_i w_i y_i=1,\\
&y_i-y_j\le0 &&(i,j)\in\Aset,\\
&y_i\ge0 &&i\in\Vset.
\end{aligned}
\end{equation}
These $y_i$ are normalized variables. They have \emph{no upper bound of one}.
For a nonempty closure $S$, the vector $\vect{y}=\ind^S/w(S)$ is feasible and has
objective $\varrho(S)$. Its coordinates can exceed one if $w(S)<1$.

\begin{theorem}[Exact ratio normalization]\label{thm:normalization}
The feasible set of \eqref{eq:normalized} is nonempty and compact. Its
optimal value is \eqref{eq:ratio}, including when this value is negative.
\end{theorem}
\begin{proof}
The full vertex set is a nonempty closure. Positivity of the weights implies
$0\le y_i\le1/w_i$, so the feasible set is bounded and closed.
For any feasible $y$, define $C_t=\{i:y_i\ge t\}$ for $t>0$.
Each $C_t$ is a closure, and the finite step function representation gives
\[
 y=\int_0^\infty\ind^{C_t}\,dt,\qquad
 1=\int_0^\infty w(C_t)\,dt,\qquad
 \vect{p}^T\vect{y}=\int_0^\infty p(C_t)\,dt.
\]
For every nonempty $C_t$, $p(C_t)\le\rho^*w(C_t)$; for an empty set both
sides vanish. Integrating proves $\vect{p}^T\vect{y}\le\rho^*$. Conversely,
$\ind^S/w(S)$ attains the ratio of any maximizing closure.
\end{proof}

\subsection{The class of linear program this is}\label{sec:class}
Three equivalent readings of \eqref{eq:normalized} fix its place among linear
programs, and each is used later.
\begin{enumerate}[label=(\alph*),leftmargin=*]
\item \textbf{One budget row over a homogeneous difference system.} The
      inequalities $y_i-y_j\le0$ are homogeneous difference constraints on a
      digraph; one dense equality row with positive coefficients is added.
\item \textbf{The base of an order cone.} The set $\{\vect{y}\ge0:y_i\le y_j\}$
      is the cone of nonnegative order-preserving vectors on the poset, whose
      extreme rays are the characteristic vectors of the nonempty closures; the
      equality slices that cone into the polytope with vertices $\ind^C/w(C)$.
      Optimizing over it therefore returns $\max_C\varrho(C)$, which is
      \cref{thm:normalization} read geometrically. It is the cone analogue of
      Stanley's order polytope, whose vertices are the characteristic vectors of
      the filters of the poset \citep{Stanley1986}.
\item \textbf{Transposed incidence plus one dense row.} Each difference
      constraint carries a $+1$ and a $-1$ in the columns of two vertex
      variables, so the matrix is the \emph{transpose} of a digraph incidence
      matrix with one dense row appended, and the right-hand side vanishes on
      the network part. This is $\mat{A}$ in \eqref{eq:standardmatrix} below.
\end{enumerate}
Reading (c) is worth keeping in mind. In the side-constrained network flow
literature the variables are arcs and the rows are vertices, and here it is the
other way round; the two formulations are linear programming duals of one
another, so the algebra transports between them, but pricing on one side
corresponds to the ratio test on the other. The same duality places
\eqref{eq:normalized} in the structural family of isotonic regression and of
total variation on a graph, which are difference systems on vertex variables as
well --- which is why \citet{Yang2026} is close prior art.

An inequality $\vect{w}^T\vect{y}\le1$ would allow $y=0$ and return zero when every
nonempty closure has negative profit. The equality in
\eqref{eq:normalized} is therefore deliberate. The outer capacitated LP
may choose nothing, but the ratio oracle must still identify the best
\emph{nonempty} closure.

\subsection{Standard form and the ratio dual}
Introduce $s_{ij}=y_j-y_i\ge0$. The standard form is
\begin{equation}\label{eq:standard}
 \vect{w}^T\vect{y}=1,\qquad y_i-y_j+s_{ij}=0\ ((i,j)\in\Aset),
 \qquad(y,s)\ge0.
\end{equation}
There are $n+m$ variables and $m+1$ independent equations: the slack
columns contain an identity matrix, and the normalization row is independent
of the arc rows. A basis has $m+1$ variables and its complement has $n-1$.
With $\mat{D}$ denoting the arc-by-vertex matrix with row $\vect{e}_i^T-\vect{e}_j^T$, write
\begin{equation}\label{eq:standardmatrix}
 \mat{A}=\begin{pmatrix}\vect{w}^T&0\\\mat{D}&\mat{I}_m\end{pmatrix},\qquad
 \vect{b}=(1,0,\ldots,0)^T,\qquad \vect{g}=(\vect{p},0).
\end{equation}
This matrix is used for independent checking, not stored densely by the
proposed efficient implementation.

The dual is
\begin{equation}\label{eq:ratio-dual}
\begin{aligned}
\min\quad &\beta\\
\text{s.t.}\quad &w_i\beta+\divg_i(\alpha)\ge p_i &&i\in\Vset,\\
&\alpha_{ij}\ge0 &&(i,j)\in\Aset,\qquad \beta\in\Rset.
\end{aligned}
\end{equation}
Here $\beta$ is free because normalization is an equality. For any closure
$S$, summing the inequalities gives
\[
 p(S)\le\beta w(S)+\sum_{i\in S}\divg_i(\alpha)\le\beta w(S).
\]
The last inequality holds because no arc leaves a closure and nonnegative
flow may enter it. Thus a feasible dual vector provides a direct certificate
of an upper bound on every closure ratio.

\section{A basis represented by a forest}\label{sec:forest}

The central combinatorial object is a pair $(\Fset,\Qset)$: an undirected
forest formed from selected precedence arcs, retaining each arc's original
orientation, and a set of root vertices. Every tree has one root except
for a single rootless tree. The theorem below establishes that this object
encodes exactly the nonsingular bases of the normalized LP. Its feasible
instances are then identified by a closure condition.

\begin{remark}[The forest is not new]
A forest carrying per-branch aggregates, with branches classified by the sign of
the aggregate they support, is the state maintained by the open-pit design
algorithm of \citet{LerchsGrossmann1965}: its \emph{normalized tree} hangs
branches from an added root, records the mass $b(T_q)$ of each branch, calls a
downward arc \emph{strong} when the mass it supports is positive, and pivots by
adding one arc and removing a branch's root arc. \citet{Hochbaum2001} sets that
algorithm out in this language and shows the tree determines a flow in linear
time. The same forest structure appears for the normalized ratio program in
\citet{Wuille2025}, and the revised-simplex algebra on rooted spanning forests
is developed in \citet{Yang2026} for a different linear program. What follows is
therefore a specialisation of the ordinary primal simplex to
\eqref{eq:standard}, not a new class of basis; the positioning document in
\path{literature/} states precisely what that leaves and what it does not.
\end{remark}

\subsection{Nonbasic slacks are the forest edges}
Let $\Nset$ denote the nonbasic variables in \eqref{eq:standard}, and define
\[
 \Fset=\{(i,j):s_{ij}\in\Nset\},\qquad
\Qset=\{i:y_i\in\Nset\}.
\]
The basic variables are recovered from the object by taking the complement:
\[
 \Bset=\{y_i:i\notin\Qset\}\ \cup\ \{s_e:e\notin\Fset\}.
\]
Thus a forest edge represents a \emph{nonbasic slack}, and a root represents
a \emph{nonbasic vertex variable}. Every remaining variable is basic, including
those whose value happens to be zero. The representation records which tight
constraints define the basis, even when additional constraints are also tight.

Setting the nonbasic variables to zero gives
\begin{equation}\label{eq:nonbasic}
 y_i=y_j\quad((i,j)\in\Fset),\qquad y_q=0\quad(q\in\Qset).
\end{equation}
Consequently, the vertex values are constant on each connected component of
$(\Vset,\Fset)$, with directions ignored only for connectivity.
The original orientation remains necessary when slack variables move.

\begin{theorem}[Forest characterization of bases]\label{thm:basis}
A set of $n-1$ nonbasic variables is the complement of a nonsingular basis
of \eqref{eq:standard} if and only if:
\begin{enumerate}[label=(\roman*)]
\item $\Fset$ is a spanning forest, with isolated vertices allowed;
\item each of its trees has at most one vertex in $\Qset$;
\item exactly one tree $T_*$ has no such vertex; every other tree has one.
\end{enumerate}
A vertex in $\Qset$ is called the root of its tree.
\end{theorem}
\begin{proof}
Eliminating all slacks leaves the normalization row and the $n-1$ equations
in \eqref{eq:nonbasic}, a square system on $y$. A cycle among forest edges
would make its edge-difference rows dependent. Likewise, two roots in one
component would be dependent modulo the path equations. Hence nonsingularity
requires an acyclic edge set and at most one root per component.
If the forest has $k$ trees, it has $n-k$ edges. Counting $n-1$ nonbasic
variables then requires $k-1$ roots, leaving exactly one tree rootless.
Conversely, the edge equations reduce $y$ to one scalar per tree. The
$k-1$ roots fix all but one scalar to zero, and normalization fixes the
remaining scalar uniquely because $w(T_*)>0$.
\end{proof}

At the resulting basic solution,
\begin{equation}\label{eq:basic-solution}
 y_i=\begin{cases}1/w(T_*),&i\in T_*,\\0,&i\notin T_*,\end{cases}
 \qquad s_{ij}=y_j-y_i,\qquad
 \rho=\vect{p}^T\vect{y}=\frac{p(T_*)}{w(T_*)}.
\end{equation}
We call $T_*$ the positive tree and all other trees zero trees.

\begin{corollary}[Feasibility and support]\label{cor:feasible}
The forest basis is primal feasible if and only if $T_*$ is a closure.
Every feasible basis has exactly one positive vertex level, and its support
is the connected closure $T_*$.
\end{corollary}
\begin{proof}
The only way a slack in \eqref{eq:basic-solution} can be negative is an arc
whose tail is in $T_*$ and whose head lies outside it. Excluding those arcs
is exactly the closure condition.
\end{proof}

This characterization concerns the normalized oracle, not every optimum of
the original capacitated LP. Moreover, a \emph{basic} variable can be zero.
An arc outside $\Fset$ has a basic slack, but that slack is zero whenever its
endpoints have the same value. These additional tight arcs account for many
degenerate pivots.

\subsection{A simple feasible initialization}
A nonempty DAG has a sink $v$, meaning that no arc leaves $v$ under our
prerequisite convention. Choose
\begin{equation}\label{eq:initial}
 \Fset=\emptyset,\qquad \Qset=\Vset\setminus\{v\},\qquad T_*=\{v\}.
\end{equation}
The basic variables are $y_v$ and all arc slacks. Set $y_v=1/w_v$ and all
other vertex variables to zero. The singleton $\{v\}$ is a closure, so this
basis is feasible. This construction works on disconnected graphs and needs
no artificial arcs or phase-I problem. For the numerical trace, choose the
smallest-index sink, vertex 1.

\subsection{Rooting and the orientation sign}
Root each zero tree at its root; root $T_*$ at any vertex $r_*$. For a tree
edge $e$, let $a$ be its parent endpoint and $b$ its child endpoint. Let
$D_e$ be the subtree rooted at $b$, including $b$, and put
\begin{equation}\label{eq:sign}
 \sigma_e=\begin{cases}
 +1,&e=(a,b)\text{ points from parent to child},\\
 -1,&e=(b,a)\text{ points from child to parent}.
 \end{cases}
\end{equation}
The slack equation implies
\begin{equation}\label{eq:parent-child}
 y_b-y_a=\sigma_e s_e.
\end{equation}
Thus a vertex value is its root value plus the signed sum of slacks on its
root-to-node path. On a general graph these signs depend on orientations,
not on path-length parity.

\section{Reduced costs and directions from subtree sums}\label{sec:directions}

We now express one revised simplex iteration using the forest. The positive
tree is $T_*$, its weight is $W_*=w(T_*)$, and the current objective is
$\rho=p(T_*)/W_*$. All formulas refer to the current basis.

\subsection{Connection with revised simplex: the two basis systems}
Write $\mat{B}=\mat{A}_{\Bset}$ and let $\vect{g}=(\vect{p},0)$ be the standard-form objective vector.
The two familiar basis systems in a primal revised simplex iteration are
\begin{align}
 \mat{B}^T\vecg{\pi}&=\vect{g}_{\Bset},
 &\bar c_q&=g_q-\mat{A}_q^T\vecg{\pi},
 &&\text{multipliers and reduced costs},\label{eq:pricing-system}\\
 Bv&=A_q,
 &d_{\Bset}&=-v,
 &&\text{column for entering variable }q.\label{eq:column-system}
\end{align}
The forest supplies explicit solutions to both systems. For the first,
subtree sums yield the reduced costs in \eqref{eq:rc-node}--\eqref{eq:rc-edge}
and the multipliers $\pi=(\rho,\alpha)$ in \eqref{eq:forest-dual}. For the
second, a signed tree or subtree movement and one normalization correction
give \eqref{eq:direction}. Thus neither stage requires a general basis
factorization or a triangular solve. The next subsections derive these
identities and explain their arithmetic costs. In this precise sense, Forest
Simplex implements the revised simplex operations through a combinatorial
representation of the basis: it computes the same multipliers, reduced costs,
and pivot directions by forest traversals and aggregate formulas.

The entering variable selects a \emph{column}; the ratio test subsequently
selects the leaving variable and hence a \emph{row}. If that pivot row itself
is required, its usual transposed basis system also has a direct forest
evaluation, given in \eqref{eq:pivot-row} below.

\subsection{The dictionary behind the formulas}
Let $a_T$ be the root variable of tree $T$. In zero trees this is the
nonbasic root $y_{q_T}$; in the positive tree it is basic. Repeated use of
\eqref{eq:parent-child} gives
\begin{equation}\label{eq:path-dictionary}
 y_i=a_{T(i)}+\sum_{e\text{ on root-to-}i\text{ path}}\sigma_e s_e.
\end{equation}
Summing with weights, each $s_e$ appears exactly at the descendants $D_e$.
Therefore
\begin{equation}\label{eq:root-dictionary}
 a_{T_*}=\frac{1-\sum_{T\ne T_*}w(T)y_{q_T}
                    -\sum_{e\in\Fset}\sigma_e w(D_e)s_e}{W_*}.
\end{equation}
Equations~\eqref{eq:path-dictionary}--\eqref{eq:root-dictionary}, followed
by $s_{ij}=y_j-y_i$ for nonforest arcs, are an explicit simplex dictionary.
They express every basic variable in terms of the $n-1$ nonbasic variables.
Only the positive root is eliminated by normalization.

Substituting into the objective yields
\begin{equation}\label{eq:objective-dictionary}
\begin{split}
 \vect{p}^T\vect{y}={}&\rho+
 \sum_{T\ne T_*}\bigl[p(T)-\rho w(T)\bigr]y_{q_T}\\
 &+\sum_{e\in\Fset}\sigma_e\bigl[p(D_e)-\rho w(D_e)\bigr]s_e.
\end{split}
\end{equation}
This identity immediately gives all reduced costs, without solving a dense
linear system.

\begin{proposition}[Forest pricing]\label{prop:pricing}
For the maximization problem, the reduced costs of the nonbasic variables are
\begin{align}
 \bar c_{y_{q_T}}&=p(T)-\rho w(T),&&T\ne T_*,\label{eq:rc-node}\\
 \bar c_{s_e}&=\sigma_e\bigl[p(D_e)-\rho w(D_e)\bigr],&&e\in\Fset.
 \label{eq:rc-edge}
\end{align}
A positive reduced cost identifies a candidate entering variable. It does
not guarantee a positive feasible step.
\end{proposition}

\paragraph{Specialized entering-variable selection.}
A postorder traversal computes all subtree profit and weight sums in $O(n)$
arithmetic operations. With these aggregates available, either formula above
costs $O(1)$ per candidate. There are exactly $n-1$ nonbasic variables, so full
pricing and Bland's entering-variable selection cost $O(n)$, independently of
the number m of original arcs. No arc outside the forest needs pricing.
Writing $P_*=p(T_*)$, signs and comparisons of reduced costs can also use
their numerators over the common positive denominator $W_*$:
\[
 W_*p(T)-P_*w(T),\qquad
 \sigma_e\bigl[W_*p(D_e)-P_*w(D_e)\bigr].
\]
These comparisons avoid division; they do not remove the need to control
overflow or rounding in a numerical implementation.

\subsection{A specialized pivot column from one component movement}
Choose a nonbasic variable $q$ and increase it from zero to $t$. Define the
unadjusted vertex displacement $h$ by
\begin{equation}\label{eq:h}
 h=\begin{cases}
 \ind^T,&q=y_{q_T}\text{ is a root},\\
 \sigma_e\ind^{D_e},&q=s_e\text{ is a forest slack}.
 \end{cases}
\end{equation}
The raw motion generally changes total weight. Restore normalization by
shifting every vertex of $T_*$ by the same amount:
\begin{equation}\label{eq:direction}
 \kappa=\frac{\sum_i w_i h_i}{W_*},\qquad
 d_i=h_i-\kappa\ind^{T_*}_i,\qquad
 d_{s_{ij}}=d_j-d_i\quad((i,j)\in\Aset).
\end{equation}
The full candidate motion is
$y_i(t)=y_i+t d_i$ and $s_{ij}(t)=s_{ij}+t d_{s_{ij}}$.

\begin{proposition}[The forest motion is the simplex direction]
\label{prop:direction}
The direction in \eqref{eq:direction} has component one at the entering
variable, zero at every other nonbasic variable, and satisfies all homogeneous
equations of \eqref{eq:standard}. Consequently, its basic components equal
$-A_{\Bset}^{-1}A_q$, and its objective derivative is $\bar c_q$.
\end{proposition}
\begin{proof}
For a root entry, $h$ is constant on each tree, and the chosen root
increases by one. All other roots remain fixed. For an edge entry, $h$
is constant on each side of the deleted tree edge and its head-minus-tail
difference on that edge is one by \eqref{eq:sign}. No other forest edge has
a nonzero head-minus-tail difference. No zero-tree root belongs to $D_e$.
The correction on $T_*$ is constant on that entire tree and affects neither
these differences nor any zero-tree root. By construction $\vect{w}^T\vect{d}=0$, and
the slack differences enforce every arc equation. Uniqueness follows from
nonsingularity of the basis. Finally,
\[
 \vect{p}^T\vect{d}=\vect{p}^T\vect{h}-\rho \vect{w}^T\vect{h},
\]
which gives \eqref{eq:rc-node} or \eqref{eq:rc-edge}.
\end{proof}

\paragraph{Computing the column without a basis solve.}
The conventional tableau column is $A_{\Bset}^{-1}A_q=-d_{\Bset}$.
Equation~\eqref{eq:direction} computes it directly from the forest, including
its sign convention. For root entry, $\kappa=w(T)/W_*$; for slack entry,
$\kappa=\sigma_e w(D_e)/W_*$. Each uses a stored aggregate.
The direction is described by the selected tree or subtree, its sign, the
positive tree, and one scalar. Once component membership is available and
the selected subtree is marked, each vertex coordinate and each slack
coordinate costs $O(1)$ to evaluate. Marking and forming all vertex coordinates
cost $O(n)$; explicitly forming all slack coordinates costs $O(m)$.
The ratio test can evaluate these coordinates as it scans the basic variables,
so storing a separate complete pivot column is optional.

\paragraph{The pivot row, explicitly.}
Suppose the ratio test chooses basic variable $z_b$, at position r in the
ordered basis. A conventional way to obtain its tableau row is to solve
$\mat{B}^T\vecg{\eta}=\vect{e}_r$ and form $a_{rq}=\vecg{\eta}^T\mat{A}_q$.
For each nonbasic candidate q, let $h^{(q)}$ and $\kappa_q$ be the quantities
in \eqref{eq:h}--\eqref{eq:direction}. The same forest formulas give directly
\begin{equation}\label{eq:pivot-row}
 a_{rq}=-d_b^{(q)}=
 \begin{cases}
 \kappa_q\ind^{T_*}_i-h_i^{(q)},&z_b=y_i,\\[2pt]
 h_i^{(q)}-h_j^{(q)}+
 \kappa_q(\ind^{T_*}_j-\ind^{T_*}_i),&z_b=s_{ij}.
 \end{cases}
\end{equation}
Indeed, $d_{\Bset}^{(q)}=-B^{-1}A_q$, so its component at position r
is the negative of the requested row coefficient. Component identifiers and
depth-first entry/exit indices allow tree and subtree membership queries in
$O(1)$ after $O(n)$ forest preprocessing. With the aggregates also available,
each coefficient above costs $O(1)$ and all $n-1$ nonbasic coefficients of a
chosen row cost $O(n)$. The basic coefficients are the corresponding identity
row. The reference primal algorithm only needs the entering column and ratio
test; \eqref{eq:pivot-row} makes explicit how to recover the row when needed.

\paragraph{Why all arcs still matter.}
Only forest arcs are needed for pricing and vertex directions. However, the
basic slack on \emph{every} nonforest arc can decrease. Those arcs must be
considered in the ratio test even when they join two different trees or are
transitively redundant. Ignoring them can produce an infeasible pivot.

\subsection{Reading a dual certificate from the forest}
Define arc numbers and vertex residuals by
\begin{align}
 \beta&=\rho,\qquad
 \alpha_e=\begin{cases}-\bar c_{s_e},&e\in\Fset,\\0,&e\notin\Fset,
 \end{cases}\label{eq:forest-dual}\\
 r_i&=w_i\rho+\divg_i(\alpha)-p_i.\label{eq:dual-residual}
\end{align}

\begin{proposition}[Optimality certificate]\label{prop:certificate}
For a forest basis, $r_i=0$ at every basic vertex variable and
$r_{q_T}=-\bar c_{y_{q_T}}$ at every root. Thus if all nonbasic reduced
costs are nonpositive, \eqref{eq:forest-dual} is feasible for
\eqref{eq:ratio-dual} and certifies the optimal ratio $\rho$.
\end{proposition}
\begin{proof}
Set $u_i=p_i-\rho w_i$. On any rooted subtree $D_e$, the only forest edge
crossing its boundary is $e$. Its net outgoing flow is
$-\sigma_e\alpha_e=u(D_e)$ by \eqref{eq:rc-edge}. Taking differences between
a subtree and its child subtrees proves $\divg_i(\alpha)=u_i$ at every
nonroot vertex. On $T_*$, $u(T_*)=0$, so the same equality holds at its root.
In a zero tree, all nonroot equalities hold and the residual at the root
is $-u(T)=\rho w(T)-p(T)$. Nonpositive reduced costs make every $\alpha_e$
and every $r_i$ nonnegative. The primal and dual objectives both equal
$\rho$, proving optimality.
\end{proof}
This certificate uses sums on the original arcs of the residual instance.
It is suitable for a verifier independent of the data structures used to
update the forest.
At every forest basis, even before optimality, the vanishing residuals at
basic vertex variables and the zero multipliers on basic slacks imply
$\mat{B}^T\vecg{\pi}=\vect{g}_{\Bset}$. Thus the construction also solves the first basis system
\eqref{eq:pricing-system}; nonpositive reduced costs are needed only to make
these multipliers dual feasible.

\section{The ratio test, forest updates, and termination}\label{sec:pivots}

\subsection{The largest admissible step}
For brevity let $\vect{z}=(\vect{y},\vect{s})$ denote the standard-form variable vector in this
section only, and let $d$ include vertex and slack components from
\eqref{eq:direction}. With entering variable $q$, every other nonbasic
variable stays zero. Nonnegativity of the basic variables gives
\begin{equation}\label{eq:step}
 t^\star=\min_{b\in\Bset:\ d_b<0}\frac{z_b}{-d_b}.
\end{equation}
Choose a leaving variable attaining this minimum, and set
\begin{equation}\label{eq:pivot}
 z'=z+t^\star d,\qquad
 \Nset'=(\Nset\setminus\{q\})\cup\{b\}.
\end{equation}
The set in \eqref{eq:step} cannot be empty: the full normalized polytope,
including its slacks, is bounded, whereas $d_q=1$ would otherwise yield
an infinite feasible ray. The new objective is
$\rho'=\rho+t^\star\bar c_q$.

If a blocking basic variable is already zero, $t^\star=0$. Such a pivot
changes the basis and the forest but not the point. It can be necessary
even when the current point already has the optimal objective value.

\subsection{All four structural cases}
Forest edges encode \emph{nonbasic} slacks. Hence an entering slack is
removed from the forest, whereas a leaving slack is added to it. Roots
obey the corresponding rule for vertex variables.
\begin{table}[htbp]
\centering\small
\begin{tabular}{llll}
\toprule
Entering & Leaving & Edge update & Root update\\
\midrule
$y_q$ & $y_b$ & $\Fset' =\Fset$ & $\Qset'=(\Qset\setminus\{q\})\cup\{b\}$\\
$y_q$ & $s_f$ & $\Fset'=\Fset\cup\{f\}$ & $\Qset'=\Qset\setminus\{q\}$\\
$s_e$ & $y_b$ & $\Fset'=\Fset\setminus\{e\}$ & $\Qset'=\Qset\cup\{b\}$\\
$s_e$ & $s_f$ & $\Fset'=(\Fset\setminus\{e\})\cup\{f\}$ & $\Qset'=\Qset$\\
\bottomrule
\end{tabular}
\caption{Updates determined by the entering and leaving variable types.}
\label{tab:updates}
\end{table}
Removing a forest edge splits a tree; adding a leaving slack joins the
appropriate components after any split. A root exchange can transfer
which component is positive. One should not guess these changes from an
improvement in the objective: the ratio test determines the valid exchange.

\begin{theorem}[Correctness of one pivot]\label{thm:pivot}
Starting from a feasible forest basis, a positive reduced-cost entry and
the full ratio test \eqref{eq:step} produce a feasible basis represented
by \cref{tab:updates}. The objective does not decrease.
\end{theorem}
\begin{proof}
\Cref{prop:direction} identifies exactly the ordinary simplex direction.
The ratio test preserves nonnegativity of every basic variable; equality
constraints and the remaining nonbasic zero values are preserved by the
direction. The pivot element is nonzero because $d_b<0$, so exchanging
the columns yields a nonsingular basis. Its nonbasic complement is exactly
\eqref{eq:pivot}, hence \cref{thm:basis} proves the forest and root
properties without additional assumptions on orientation. The objective
change is $t^\star\bar c_q\ge0$.
\end{proof}

Rebuilding the new components and evaluating \eqref{eq:basic-solution}
must reproduce $z'$. This is a particularly useful implementation check.
A separate heuristic ``tree feasibility'' test cannot replace the complete
ratio test, which includes all decreasing vertex variables and all decreasing
basic slacks.

\subsection{A finite reference algorithm}
Fix a total ordering of all vertex and slack variables. For the reference
version use Bland's rule: the smallest-index nonbasic variable with
$\bar c_q>0$ enters; among all minimum-ratio blockers the smallest-index
basic variable leaves. Both parts of this rule matter for the standard
finiteness guarantee \citep{Bland1977}.

\begin{algorithm}[htbp]
\caption{\FSC{} ratio oracle, exact-arithmetic reference}
\label{alg:forest}
\begin{algorithmic}[1]
\Require Nonempty DAG, profits $p$, positive weights $w$, fixed variable order
\Ensure Optimal ratio $\rho$, an optimal connected closure $T_*$, dual certificate
\State Choose a sink $v$; set $\Fset\gets\emptyset$, $\Qset\gets\Vset\setminus\{v\}$
\Loop
  \State Recover trees, roots and $T_*$; root the trees
  \State Compute subtree $p,w$ sums, basic values, and $\rho$ from \eqref{eq:basic-solution}
  \State Price all nonbasic variables using \eqref{eq:rc-node}--\eqref{eq:rc-edge}
  \If{all reduced costs are nonpositive}
    \State \Return $\rho,T_*$ and the certificate \eqref{eq:forest-dual}
  \EndIf
  \State Select the smallest-index improving nonbasic variable $q$
  \State Compute $h$, vertex directions and all slack directions by \eqref{eq:direction}
  \State Compute $t^\star$ and the smallest-index tied leaving variable $b$ by \eqref{eq:step}
  \State Update $\Fset,\Qset$ using \cref{tab:updates}
\EndLoop
\end{algorithmic}
\end{algorithm}

\begin{theorem}[Finite exact-arithmetic oracle]\label{thm:finite}
\Cref{alg:forest} terminates at an optimum of \eqref{eq:ratio}.
\end{theorem}
\begin{proof}
Initialization is feasible by \eqref{eq:initial}. Each iteration is the
ordinary primal simplex algorithm with Bland's entering and leaving rules,
by \cref{prop:direction,thm:pivot}. The normalized polytope is nonempty
and bounded. Bland's finite pivoting theorem therefore applies. At
termination \cref{prop:certificate} provides an optimality certificate,
and \cref{thm:normalization} identifies its value with the closure ratio.
\end{proof}
This is a finite-pivot theorem in exact arithmetic, not a polynomial bound
on the number of pivots, and not a guarantee for a floating-point variant
with a different pricing rule. Such variants require their own safeguards.

\section{A complete execution on the running instance}\label{sec:trace}

\subsection{Variable order and the eight pivots}
Order $y_1,\ldots,y_8$ first, followed by the nine slacks in the arc order
\eqref{eq:arcs}. Initialize with sink 1: $y_1=1$, all other vertex variables
zero, $\Fset=\emptyset$, and roots $2,\ldots,8$. The initial ratio is 1.
The exact reference execution is shown in \cref{tab:trace}.

\begin{table}[htbp]
\renewcommand{\arraystretch}{1.4}
\centering\input{../build/tutorial/generated/trace_table.tex}
\caption{Complete Bland-rule trace for the first ratio oracle. Each row is
checked against a separately assembled rational basis system.}\label{tab:trace}
\end{table}

Only pivots 4 and 7 have a positive step. The first three pivots make
dependencies tight before vertices 5 and 6 can increase. Pivots 5 and 6
also reorganize a zero region. Pivot 8 is needed to remove a positive
reduced cost from a degenerate optimal basis.

\subsection{A first zero step: a prerequisite blocks motion}
At the initial basis, $y_5$ has reduced cost $6-1\cdot2=4$.
Its raw displacement is $h=\ind^{\{5\}}$, and normalization gives
$d_5=1$, $d_1=-2$, with all other vertex directions zero. Consequently
\[
 d_{s_{5,1}}=-3,\qquad d_{s_{5,2}}=-1,\qquad d_{s_{5,6}}=-1.
\]
Although $s_{5,1}=1$, both $s_{5,2}$ and $s_{5,6}$ are zero. The step is
zero, and the fixed variable order makes $s_{5,2}$ leave. The new forest
contains $(5,2)$, its zero tree $\{2,5\}$ is rooted at 2, and the primal
point is unchanged. A positive reduced cost alone was insufficient to add
vertex 5: its prerequisites had to be handled first.

\subsection{The principal split: pivot 7}
Immediately before pivot 7, the positive tree and zero tree are
\[
 T_*=\{1,2,3,5,6\},\qquad T_0=\{4,7,8\},
\]
with $w(T_*)=6$, $p(T_*)=10$, and $\rho=5/3$. Their forest edges are
\[
 \Fset=\{(5,1),(5,2),(5,6),(6,3),(7,4),(8,7)\},
 \qquad \Qset=\{4\}.
\]
Root $T_*$ at 1. For $e=(5,6)$, parent 5 points to child 6, so
$\sigma_e=+1$ and $D_e=\{3,6\}$. Its aggregates are $p(D_e)=4$,
$w(D_e)=2$. Therefore
\[
 \bar c_{s_{5,6}}=4-\frac53\,2=\frac23>0,\qquad
 \kappa=\frac{2}{6}=\frac13.
\]
The vertex direction, in vertex order, is
\begin{equation}\label{eq:example-direction}
 d_y=\left(-\tfrac13,-\tfrac13,\tfrac23,0,-\tfrac13,\tfrac23,0,0\right).
\end{equation}
Each positive vertex initially has value $1/6$. The decreasing basic vertex
variables $y_1,y_2,y_5$ all give ratio $1/2$. So does the basic slack
$s_{8,5}=1/6$, whose direction is $d_5-d_8=-1/3$.
The smallest-index tied blocker is $y_1$, hence $t^\star=1/2$.
The update removes $(5,6)$ from $\Fset$ and adds root 1.

The new positive tree is $\{3,6\}$, with value $1/2$ at both vertices.
The vertices $\{1,2,5\}$ form a zero tree rooted at 1. The ratio increases to
\[
 \frac53+\frac12\cdot\frac23=2.
\]
\begin{figure}[htbp]
\centering
\begin{minipage}{.49\textwidth}\centering
\resizebox{\linewidth}{!}{\input{../build/tutorial/generated/forest_before.tex}}
\small Before pivot 7: positive tree $\{1,2,3,5,6\}$.
\end{minipage}\hfill
\begin{minipage}{.49\textwidth}\centering
\resizebox{\linewidth}{!}{\input{../build/tutorial/generated/forest_after.tex}}
\small After pivot 7: positive tree $\{3,6\}$.
\end{minipage}
\caption{The nonbasic-slack forest in solid blue, other arcs dashed.
Positive-tree vertices are shaded blue; double borders mark roots. The
right-hand basis still requires the degenerate pivot 8.}\label{fig:split}
\end{figure}

\subsection{An optimal point need not yet have an optimal basis}
After pivot 7, consider edge $(5,1)$ in the zero tree $\{1,2,5\}$,
rooted at 1. It points from child 5 to parent 1, so $\sigma=-1$ and
$D=\{2,5\}$. Pricing gives
\[
 \bar c_{s_{5,1}}=-\bigl[5-2\cdot3\bigr]=1>0.
\]
Its raw motion decreases $y_2$ and $y_5$, which are already zero. The
ratio test returns zero and $y_2$ leaves. The objective remains 2, but
the forest now has roots $1,2,4$ and edges
\[
 \Fset=\{(5,2),(6,3),(7,4),(8,7)\}.
\]
All nonbasic reduced costs are now nonpositive. The ratio certificate is
\begin{equation}\label{eq:ratio-example-cert}
 \beta=2,\qquad
 \alpha_{5,2}=2,\quad\alpha_{6,3}=3,\quad
 \alpha_{7,4}=1,\quad\alpha_{8,7}=1,
\end{equation}
with every other arc multiplier zero. Its vertex residuals are 1, 1 and 4
at roots 1, 2 and 4, respectively, and zero at all other vertices.
It proves that no nonempty closure has ratio greater than 2.

This final pivot is a useful review case: stopping solely because the
objective did not change, or because the support already looks correct,
would bypass the basis optimality test.

\section{From the forest oracle to the capacitated LP}\label{sec:outer}

The \FS{} returns an optimal \emph{connected support} of the normalized
oracle. We now explain how repeated calls solve the original LP, how equal
ratios are handled, and why this outer construction is valid on general
precedence graphs.

\subsection{Residual extraction and decreasing ratios}
Start with $M=\emptyset$. Call the oracle on the residual graph, obtain its
positive tree $B$, append $B$ to the sequence and replace $M$ by $M\cup B$.
The set $B$ is closed on the residual graph, so every accumulated prefix
remains closed on the original graph. Each call removes at least one vertex.
Thus at most $n$ calls partition the vertex set, even if some ratios are zero
or negative. Denote the raw blocks by $B_1,\ldots,B_\ell$ and their ratios
by $\rho_1,\ldots,\rho_\ell$.

\begin{lemma}[Nonincreasing raw ratios]\label{lem:ratios}
The residual extraction procedure satisfies $\rho_{r+1}\le\rho_r$.
\end{lemma}
\begin{proof}
The union $B_r\cup B_{r+1}$ is closed in the residual graph before call r.
Its ratio is the weighted average of $\rho_r$ and $\rho_{r+1}$, with strictly
positive weights. A value $\rho_{r+1}>\rho_r$ would make this union better
than the closure returned at call r, contradicting its certified optimality.
\end{proof}

\subsection{Why a single basis may miss a canonical block}
For two independent vertices with $p_1=w_1=p_2=w_2=1$, either singleton is a
valid optimal positive tree. However, the maximum-weight optimal closure
is $\{1,2\}$. A single basic support need not equal a canonical block.
Merge consecutive raw blocks with the same ratio. Write the resulting blocks
as $\Kset_1,\ldots,\Kset_k$, their closed prefixes as $\mathcal M_r$, and
their strictly decreasing ratios as $\lambda_r$.

The following proof establishes both the LP connection and the maximal tie
convention directly. It is consistent with the classical parametric closure
interpretation reviewed in \citet{DoseFuriniLocatelli2026} and
\citet{GalloEtAl1989}.

\begin{lemma}[Blockwise closure bound]\label{lem:blockbound}
For any original closure $C$ and every raw block $B_r$,
\[
 p(C\cap B_r)\le\rho_r w(C\cap B_r).
\]
Consequently, for every real $\lambda$,
\begin{equation}\label{eq:parametric-sum}
 u(\lambda):=\max_{C\in\Cset}\{p(C)-\lambda w(C)\}
 =\sum_{r=1}^\ell w(B_r)\max\{\rho_r-\lambda,0\}.
\end{equation}
\end{lemma}
\begin{proof}
Both $C$ restricted to the residual graph and $B_r$ are residual closures,
and their intersection is one too. If it is nonempty, the oracle certificate
at call r bounds its ratio by $\rho_r$; otherwise both sums vanish. Hence
\[
 p(C)-\lambda w(C)
 \le\sum_r(\rho_r-\lambda)w(C\cap B_r)
 \le\sum_r w(B_r)\max\{\rho_r-\lambda,0\}.
\]
The union of all raw blocks with $\rho_r>\lambda$ is a prefix by
\cref{lem:ratios}; it is a closure and attains this bound.
\end{proof}

\begin{corollary}[Canonical merging]\label{cor:canonical}
Merging all consecutive equal-ratio raw blocks produces exactly the
maximum-weight tie convention of the supplied manuscript.
\end{corollary}
\begin{proof}
At any $\lambda$, equality in the second inequality of the preceding proof
forces an optimal closure to contain every vertex of a block with
$\rho_r>\lambda$ and no vertex of a block with $\rho_r<\lambda$; positivity
of each vertex weight is used here. The prefix including every block with
$\rho_r=\lambda$ is feasible, attains the same value and contains every
optimal closure. It is the unique maximum-weight one. At each distinct
ratio, its increment is precisely the group merged by the procedure.
\end{proof}

This proof also explains why canonical merging is not an optional heuristic.
Equal-ratio blocks may be disconnected. A solver asked for a complete
canonical block must finish collecting its ties, even if a requested
capacity has already been reached.

\subsection{The canonical LP solution}
Let $q$ be the last index for which $\lambda_q>0$, with $q=0$ when no such
index exists. If
$w(\mathcal M_{h-1})\le c<w(\mathcal M_h)$ for some $h\le q$, set
\begin{equation}\label{eq:theta}
 \theta=\frac{c-w(\mathcal M_{h-1})}{w(\Kset_h)},\qquad
 \vect{x}^c=\ind^{\mathcal M_{h-1}}+\theta\ind^{\Kset_h}.
\end{equation}
If $c\ge w(\mathcal M_q)$, set $\vect{x}^c=\ind^{\mathcal M_q}$, including the
empty solution when $q=0$.

\begin{theorem}[Forest-based solution of the original LP]\label{thm:outer}
The vector $\vect{x}^c$ above is optimal for \eqref{eq:lp}. Within a positive
block interval,
\begin{equation}\label{eq:value}
 z(c)=p(\mathcal M_{h-1})+
              \lambda_h\bigl(c-w(\mathcal M_{h-1})\bigr).
\end{equation}
For $c\ge w(\mathcal M_q)$, $z(c)=p(\mathcal M_q)$.
\end{theorem}
\begin{proof}
The vector in \eqref{eq:theta} is a convex combination of two closure
indicators, so it respects precedences and box constraints; it uses capacity
c. For any feasible x, its threshold sets $C_t=\{i:x_i\ge t\}$,
$0<t\le1$, are closures. Thus for any $\lambda\ge0$,
\[
 \vect{p}^T\vect{x}-\lambda \vect{w}^T\vect{x}
   =\int_0^1[p(C_t)-\lambda w(C_t)]\,dt\le u(\lambda),
 \qquad \vect{p}^T\vect{x}\le u(\lambda)+\lambda c.
\]
Choose $\lambda=\lambda_h$ and use \eqref{eq:parametric-sum}. Both adjacent
prefixes attain $u(\lambda_h)$, so $\vect{x}^c$ attains the bound. For abundant
capacity choose $\lambda=0$; the positive prefix attains $u(0)$.
\end{proof}

The theorem guarantees a canonical optimum with at most one fractional
level; it does not claim that every optimal point has that form. A zero-ratio
block may be included without changing the objective if capacity allows.
Negative-ratio blocks must not be added merely to fill unused capacity.

\begin{algorithm}[htbp]
\caption{Block extraction using only \FS{} oracles}
\label{alg:outer}
\begin{algorithmic}[1]
\Require General precedence instance after SCC contraction
\Ensure Complete canonical sequence, or its positive part as requested
\State $R\gets\Vset$; sequence $\gets$ empty
\While{$R\ne\emptyset$}
  \State $(\rho,B,\text{certificate})\gets\FS(R)$ using \cref{alg:forest}
  \If{only the positive path is requested and $\rho\le0$}
    \State \textbf{break}
  \EndIf
  \If{last block has ratio $\rho$}
    \State Merge B into the last block
  \Else
    \State Append B with ratio $\rho$
  \EndIf
  \State $R\gets R\setminus B$
\EndWhile
\State Use \eqref{eq:theta} to answer a prescribed capacity, if requested
\end{algorithmic}
\end{algorithm}

\subsection{Finishing the eight-node example}
Successive calls and canonical merging give \cref{tab:blocks}.
\begin{table}[htbp]
\centering\input{../build/tutorial/generated/blocks_table.tex}
\caption{Canonical blocks produced by residual \FS{} calls,
checked against exhaustive closure enumeration.}\label{tab:blocks}
\end{table}
At $c=4$, the first block has weight 2, so half the second block is selected:
\[
 \theta=\frac{4-2}{4}=\frac12,\qquad
 x^c=(\tfrac12,\tfrac12,1,0,\tfrac12,1,0,0),\qquad z(4)=4+\tfrac12\,6=7.
\]
A full dual certificate for \eqref{eq:lpdual} is
\begin{align*}
 \lambda&=\tfrac32,\quad\mu_6=1,\quad\mu_i=0\ (i\ne6),\\
 \alpha_{5,1}&=\tfrac12,\quad\alpha_{5,2}=\tfrac52,\quad
 \alpha_{6,3}=\tfrac52,\quad\alpha_{7,4}=2,\quad\alpha_{8,7}=\tfrac32,
\end{align*}
with all remaining arc multipliers zero. Its objective is
$4\cdot(3/2)+1=7$. Direct substitution verifies feasibility and complementary
slackness. This certificate differs from the normalized certificate
\eqref{eq:ratio-example-cert}: the original LP has capacity and upper-bound
prices, while the ratio oracle has only a free normalization price.

\section{Efficient implementation: what the derivation permits}\label{sec:implementation}

The exact reference algorithm is fully specified above. This section describes
the implementable baseline, and then the three refinements that measurement
showed to matter. The baseline is not a straw man: it is the configuration the
solver still offers as its verifiable reference, and every refinement below is a
separate option so that its effect can be measured in isolation.

\subsection{A linear-memory reference organization}
Store the original directed graph in contiguous endpoint and adjacency arrays,
with incoming and outgoing access. Store the forest separately through parent,
parent-edge, root and child information. Maintain the nonbasic vertex/slack
sets and a mapping between variable IDs and their positions in basic lists.
Traverse trees iteratively, so a long chain cannot exhaust the call stack.

A postorder traversal computes $p(D_e)$ and $w(D_e)$ for every forest edge.
Component totals provide root reduced costs. Only $n-1$ nonbasic variables
need pricing. After selecting an entry, mark its subtree or tree, form
\eqref{eq:direction}, and scan all arcs for decreasing basic slacks.
Rebuild the forest metadata after the exchange. This deliberately simple
version provides a reliable reference before local updates are introduced.

\subsection{Specialized operations in one iteration}
The forest supplies specialized operations for entering-variable selection,
the pivot column, the leaving-variable test, and the basis exchange.
Table~\ref{tab:operation-costs} separates these costs for the full-rebuild
reference organization. In particular, the $O(n)$ pricing bound is independent
of graph density; the complete direction and ratio test also inspect arcs.

\begin{table}[htbp]
\centering\small
\renewcommand{\arraystretch}{1.15}
\begin{tabular}{p{.29\linewidth}p{.49\linewidth}l}
\toprule
Operation & Forest calculation & Cost\\
\midrule
Subtree aggregates & Postorder sums of profits and weights & $O(n)$\\
One reduced cost & Signed aggregate difference, sums available & $O(1)$\\
Full pricing and entry & Scan the $n-1$ nonbasic candidates & $O(n)$\\
Pivot direction on vertices & Mark a tree/subtree; normalize its motion & $O(n)$\\
Pivot direction on slacks & Head-minus-tail vertex directions & $O(m)$\\
Ratio test and exit & Scan decreasing basic vertex/slack variables & $O(n+m)$\\
Chosen pivot row, if needed & Evaluate its $n-1$ nonbasic coefficients from forest membership and aggregates & $O(n)$\\
Basis exchange and rebuild & Exchange edge/root membership; rebuild metadata & $O(n+m)$\\
\bottomrule
\end{tabular}
\caption{Arithmetic costs with the current forest available and reusable
arrays. Constant-time coordinate evaluation assumes the required sums and
membership information have been computed.}\label{tab:operation-costs}
\end{table}

\begin{proposition}[Arithmetic cost of the full-rebuild variant]
\label{prop:cost}
Using adjacency lists and reusable arrays, one full rebuild, complete pricing,
direction computation, full ratio test and exchange can be performed in
$O(n+m)$ arithmetic operations and $O(n+m)$ storage. An oracle performing
K pivots costs $O((K+1)(n+m))$ arithmetic operations.
\end{proposition}
\begin{proof}
Forest construction and traversals touch $O(n+m)$ stored entries. Subtree
sums and pricing touch $O(n)$ entries, since a forest has fewer than n edges.
Constructing vertex directions is linear in n, and evaluating all slack
directions and blockers is linear in m. Updating membership and rebuilding
the adjacency information is linear as well. Arrays store only vertices, arcs,
membership, aggregates and temporary vectors.
\end{proof}
This is an arithmetic-operation bound, not a bit-complexity bound for growing
rational numbers. Nor does it bound K polynomially. The exact checker
deliberately builds dense matrices and therefore does \emph{not} implement
the memory bound in \cref{prop:cost}.
For residual calls indexed by r, the analogous total is
\[
 O\!\left(\sum_r(K_r+1)(n_r+m_r)\right),
\]
when residual graphs are rebuilt linearly at each call. The number of calls
is at most the original n; no parametric-flow complexity bound is being
transferred to the \FS{}.

\subsection{Variants to implement and tune}
\begin{table}[htbp]
\centering\small
\begin{tabular}{p{.22\linewidth}p{.32\linewidth}p{.35\linewidth}}
\toprule
Parameter family & Candidate variants & Invariant to preserve\\
\midrule
Pricing & Full scan; blocked scan; candidate list & A full check before claiming no improving entry\\
Entering rule & Bland; first or best improving & A documented anti-cycling policy for each alternative\\
Forest update & Full; local; adaptive rebuild & Same mathematical basis and direction as a full reconstruction\\
Rebuild trigger & Pivot interval; changed-tree fraction; residual threshold & Recover accuracy without silently changing the problem\\
Initial basis & Sink singleton; closure of a promising seed; best undirected component & Primal feasibility: the positive component must be closed in the residual\\
Ratio test & Full scan; restricted to boundary-crossing arcs; stop at the first zero step & Same step, and the same leaving variable whenever Bland's rule is active\\
Residual start & Fresh sink basis; validated warm start & Primal feasibility on the new residual graph\\
Vertex order & Input; topological; seeded & Identical certified objective and canonical decomposition\\
\bottomrule
\end{tabular}
\caption{Tuning choices of the implementation. Every row is an explicit
parameter, so an ablation can switch off one refinement at a time.}
\end{table}
Reducing the number of visited vertices or tested candidates is useful only if
it reduces total runtime at the required accuracy. The profiler separates
pricing, directions, ratio tests, exchange, rebuilds and canonicalization. The
measured ranking of these choices is reported in the experimental report, not
here; the one result worth repeating is that the entering rule barely moves the
outcome while the treatment of degeneracy moves it by factors.

\paragraph{A caution about termination.}
First-improving or best-improving pricing is not automatically Bland's rule.
Switching to a finite rule permanently after a documented trigger, or using
a separately justified lexicographic implementation, are possible safeguards.
An informal fallback that repeatedly switches back to an unproved rule does
not inherit the finite-pivot theorem without additional reasoning. The
reasoning is supplied in \cref{prop:episodes} below, because a permanent switch
turns out to be expensive.

\subsection{Degeneracy is the dominant cost}\label{sec:degeneracy-cost}

In a basis the vertex values are constant on the positive component and zero on
every other tree. A step is therefore blocked at length zero whenever the
blocking variable is one of the many already at zero: a vertex of the entering
subtree whose direction component is negative, or a precedence slack between two
vertices that are both at zero. Degeneracy is not incidental here; it is what the
precedence structure produces.

The effect is large enough to reorganize the implementation around it. On
applied instances more than $99\%$ of the pivots of an oracle have zero step; on
a family with no arcs at all, where no slack can block, the share is zero. Three
consequences follow, and each was measured.

\paragraph{The anti-cycling fallback must not become the operating rule.}
A permanent switch to Bland's rule after a run of degenerate pivots engages
immediately and never releases, because the run never ends: every entering rule
then collapses onto Bland within the first oracle. Releasing the fallback at the
first pivot with a positive step, and placing the trigger far away, keeps the
improving rule without giving up finiteness.

\begin{proposition}[Finiteness with a released fallback]\label{prop:episodes}
Suppose the solver uses an arbitrary entering rule, switches to Bland's rule
after a fixed number of consecutive zero-step pivots, and returns to the
original rule at the first pivot with a positive step. Then the oracle performs
finitely many pivots.
\end{proposition}
\begin{proof}
Call an \emph{episode} a maximal run of iterations during which Bland's rule is
active. Within an episode every pivot has zero step, so the basic solution never
changes, and the iterations of the episode are exactly the iterations of Bland's
rule applied to the linear program from the basis at which the episode began;
by the finite-pivot theorem for Bland's rule the episode cannot cycle and must
end. An episode ends only with a pivot of positive step, which strictly
increases the objective of the normalized subproblem. Since there are finitely
many bases, hence finitely many objective values, only finitely many episodes
can occur. Between episodes only finitely many pivots are possible for the same
reason.
\end{proof}

\paragraph{The ratio test need not scan every arc.} By
\eqref{eq:direction} the direction is $+\sigma$ on the entering subtree,
$-\kappa$ on the positive component and zero elsewhere. An arc whose two
endpoints lie outside both sets therefore has zero direction difference and
cannot block. It suffices to scan the arcs incident to the entering subtree,
together with the arcs that cross the boundary of the positive component; and
since the positive component is closed, such an arc has its tail outside it, so
scanning the outgoing arcs of the vertices outside the positive component completes
the cover.

\paragraph{A zero step is the minimum, so the scan may stop at the first one.}
No step is negative, so the first blocking variable attaining zero already
determines the step. The shortcut must be suspended while Bland's rule is
active, since that rule requires the smallest index among ties. This is the
technique the group's 2023 prototype implemented as an explicit feasibility
check during the forest traversal, and it is the single largest saving in the
ratio test: on one oracle of a mining instance with $1060$ vertices the arcs
examined fall from $1\,110\,883$ to $184$, at an unchanged number of pivots.

\subsection{The dual side}\label{sec:dual}

Primal degeneracy and dual degeneracy are different defects, and on this problem
they are wildly asymmetric. Primal degeneracy is a per-pivot phenomenon: more
than $99\%$ of the pivots of an oracle have zero step. The only source of dual
degeneracy is a tie between ratios, which is a per-block phenomenon: on a
generated instance of two thousand vertices, ties affect $123$ of $267$ blocks. Even
if every tie produced one degenerate dual pivot, that is a hundred against half a
million. The dual is therefore the natural method here, which is also what the
classical guidance says: use the dual when the primal is heavily degenerate.

Two structural facts make the dual unusually simple on this problem.

\begin{proposition}[Primal infeasibility is combinatorial]\label{prop:dual-shape}
In any basis of the normalised subproblem, the vertex variables satisfy $y\ge 0$.
Every primal infeasibility is a precedence slack $s_{ij}=y_j-y_i<0$, that is an
arc leaving the positive component $P$; and the basis is primal feasible exactly
when $P$ is closed.
\end{proposition}
\begin{proof}
The basic solution takes the value $1/w(P)$ on $P$ and zero on every other tree,
so $y\ge0$ always. For an arc $(i,j)$, $s_{ij}=y_j-y_i$ is negative only when
$y_i>y_j$, which forces $i\in P$ and $j\notin P$. Conversely if $P$ is closed no
arc leaves it, all such slacks vanish or are positive, and the basis is feasible.
\end{proof}

Choosing the leaving variable is therefore a search for an unsatisfied
prerequisite, with no numerical scan, and the stopping rule is that $P$ is
closed. The starting basis is dual feasible by construction: make every vertex its
own tree and let $P$ be the vertex of maximum ratio; the reduced cost of any other
vertex $k$ is then $w_k(p_k/w_k-\rho)\le0$, and there are no forest edges yet. It
costs $O(n)$, and it is essentially forced: a singleton start is dual feasible
only for the maximum-ratio vertex.

The row of the leaving slack is as local as the primal direction. Entering a
nonbasic $q$ by one unit moves every vertex by
$\sigma_q\ind[u\in S_q]-\kappa_q\ind[u\in P]$, so the row of the slack of
$(t,h)$ is nonzero only for the root of each endpoint's tree, for the forest
edges on the two paths to their roots, and through the rank-one term $\kappa_q$
when exactly one endpoint lies in $P$. It costs $O(n+\mathrm{depth})$.

\paragraph{What the measurement says.} The dual removes the degeneracy it was
predicted to remove --- the share of zero-step pivots falls from above $99\%$ to
between $0\%$ and $21\%$ --- but the gain in time is between $1.1$ and $2.7$,
not the two orders of magnitude a naive extrapolation would suggest. The reason
is worth recording: the primal performs about half a million pivots of which only
some fifteen hundred move the objective, while the dual performs three hundred
thousand of which most do. It has traded many null steps for many tiny ones.
Choosing the leaving arc by structural criteria rather than by smallest index
--- most attractive head, largest tree, most negative reduced cost --- was tried
and does not help, because every violated slack has the same value $-1/w(P)$ and
there is nothing for a most-infeasible rule to discriminate.

\paragraph{What remains.} With those three changes the ratio test is no longer
the bottleneck. Pricing and forest rebuilding are: on a layered DAG with dense
connections between levels they account for roughly a third and a half of the
solver's time respectively, because both are $\Theta$ of the touched part of the
graph at every pivot while the pivot itself changes very little. Incremental
pricing, updated only on the components a degenerate pivot touches, and updates
confined to the modified paths, are the natural next steps; neither is
implemented here.

\subsection{Numerical checks and output contracts}
Floating-point implementations must use declared absolute and relative
tolerances, inspect primal and dual residuals in the original units, and
reconstruct values when incremental updates drift. Near-equal ratios require
particular care: approximate merging must not be mistaken for an exact tie.
On integer/rational review cases, exact comparisons provide a useful check.

The intended outputs are: ratio and connected support for an oracle;
canonical blocks and breakpoints for a path computation; a feasible x and
validated bound/certificate for a capacity solve. Time, iteration or numerical
limits must produce distinct states rather than an ``optimal'' label.
Optional internal checks can vary with configuration, but final verification
remains mandatory for results used in a scientific comparison.

Once the positive path is known, finding a capacity's split block costs
$O(\log k)$ with cumulative weights. Returning only the objective or a compact
block description can be correspondingly cheap. Materializing all n vertex
coordinates still costs $O(n)$. The many-capacity experiment must include the
initial path computation and distinguish these output requirements.

\section{What has been checked, and how to review this draft}\label{sec:verification}

\subsection{Independent exact checks}
The companion script \texttt{tools/verify/verify\_tutorial.py} uses Python's
rational arithmetic and does not depend on a commercial LP solver. It
implements the formulas in this draft as a pedagogical checker, not as the
planned high-performance solver. At each visited basis it independently
assembles $A_{\Bset}$ and solves
\[
 A_{\Bset}z_{\Bset}=b,\qquad
 A_{\Bset}d_{\Bset}=-A_q
\]
by rational Gaussian elimination. It compares every nonbasic candidate's
direction and reduced cost with the forest formulas, checks the exchanged
basis and verifies the final ratio dual.
It also solves the revised simplex multiplier system
$\mat{A}_{\Bset}^T\vecg{\pi}=\vect{g}_{\Bset}$ independently at every visited basis and compares
its solution with the forest multipliers. At each pivot, it solves
$\mat{A}_{\Bset}^T\vecg{\eta}=\vect{e}_r$ for the leaving position and checks every
nonbasic coefficient of the resulting pivot row against the forest direction.

Separately, exhaustive enumeration lists all closures on the small instance,
computes exact maximum ratios and maximal ties, and checks the complete
canonical sequence. For the LP value it enumerates capacity intersections
of pairs of closure points in the weight--profit plane; this independent
oracle does not restrict those pairs to the sequence found by the \FS{}.

\input{../build/tutorial/generated/verification_stats.tex}
The supplied deterministic check suite contains \VerifiedInstances{} instances:
the eight-node example, four explicit edge cases, and 32 seeded small random
DAGs with profits of either sign and positive integer weights. It checks
\VerifiedBases{} bases and \VerifiedDirections{} candidate directions.
The additional revised simplex comparisons check \VerifiedMultipliers{}
multiplier systems and \VerifiedPivotRows{} pivot rows.
The eight-node graph has 36 closures. Its first oracle takes eight pivots,
six at zero step, and its capacity-4 LP value is 7.
Both numerical certificates displayed in the main text are verified by
the checker or its returned forest certificate.

These counts describe a finite review suite, not comprehensive testing of
an optimized implementation, not a statistical experiment, and not evidence
of a runtime advantage. The seed and hashes of the fixture and verification
script are saved with the generated results.

\subsection{Rebuilding the tables and PDF}
From the project root, run
\begin{quote}\ttfamily
python3 tools/build\_tutorial.py
\end{quote}
The command reruns the exact checks, regenerates numerical tables and forest
figures, and compiles this document. The single output PDF is
\begin{quote}\ttfamily\small
build/tutorial/forest\_simplex\_tutorial.pdf
\end{quote}
The generated verification record is
\begin{quote}\ttfamily\small
build/tutorial/generated/verification.json
\end{quote}
All intermediate TeX files remain under \texttt{build/tutorial}. The source
instance, shared notation and bibliography each have one authoritative file.
No original project source is altered or required during this rebuild.

\subsection{Suggested order for mathematical review}
The following checkpoints identify concrete places where an error would
affect the proposed general algorithm:
\begin{enumerate}
\item Check the prerequisite direction in \eqref{eq:lpprec} and the
parent--child sign in \eqref{eq:sign}.
\item Check the converse as well as the forward implication of the basis
characterization, \cref{thm:basis}.
\item Check that the normalization correction in \eqref{eq:direction}
preserves every other nonbasic variable.
\item Check the sign of \eqref{eq:rc-edge}, especially when a child points
to its parent in a zero tree.
\item Check the inclusion of decreasing nonforest slacks in the ratio test.
\item Follow pivots 7 and 8, including the unchanged objective at pivot 8.
\item Check the sign convention in the forest dual certificate and distinguish
it from the capacitated LP dual.
\item Check that merging raw ties yields the maximal, possibly disconnected,
canonical block.
\item Check the stop at nonpositive ratios for the capacity-inequality LP.
\end{enumerate}

\subsection{Work remaining after this review version}
The exact-reference derivation above is accompanied by proofs and finite
checks. The following tasks remain open as project deliverables:
\begin{itemize}
\item external mathematical review and comparison with existing forest,
network-simplex and closure algorithms to determine the actual novelty;
\item a C++ implementation with linear-memory data structures;
\item verified local updates, alternate pricing and warm starts;
\item numerical stress tests beyond the exact small-DAG checker;
\item competitive backend comparisons, parameter tuning and reproducible
runtime/memory experiments;
\item extension beyond the declared positive-weight, single-capacity model.
\end{itemize}
In particular, this document supplies no invented benchmark results and makes
no claim that the eventual optimized solver is faster than parametric flow
or a general LP solver.

\appendix
\section{Closure and min-cut as comparison tools}\label{app:flow}

This appendix records the alternative closure oracle only to locate the
\FSC{} within the existing framework. Maximum closure admits
a minimum-cut representation \citep{Picard1976}; parametric maximum flow
and ratio-closure applications are treated by \citet{GalloEtAl1989}. The
graph-based predecessor of both is \citet{LerchsGrossmann1965}, whose
normalized tree is a forest with branch aggregates rather than a flow;
\citet{Hochbaum2001} recovers a flow from that tree in linear time and gives a
parametric variant with the cost of a single run, which makes it a comparison
tool on the ratio task as well.
These are comparison tools, not subroutines required by
\cref{alg:forest,alg:outer}.

For a fixed real parameter $\lambda$, assign vertex value
$a_i=p_i-\lambda w_i$. Add source s and sink t, capacities
\[
 c_{si}=\max\{a_i,0\},\qquad
 c_{it}=\max\{-a_i,0\},
\]
and a precedence-arc capacity $H>\sum_i\max\{a_i,0\}$.
Every minimum cut has a source-side vertex set C that is a closure, and its
capacity is
\[
 \sum_i\max\{a_i,0\}-[p(C)-\lambda w(C)].
\]
Thus minimizing the cut maximizes the penalized closure objective.
For a fixed $\lambda$, this construction works with either sign of the
parameter. A linear parametric implementation can use a different
normalization on a specified interval, as explained in the supplied manuscript.

At a breakpoint, the closure maximizing the penalized objective need not be
unique. In a residual network of a maximum flow, source-reachable vertices
give the \emph{minimal} source side. The complement of vertices that can
reach t gives the \emph{maximal} source side. A canonical comparison must use
the latter convention. For a ratio root, the empty closure can also tie at
value zero, so an implementation must preserve or recover a nonempty ratio
maximizer.

The planned experiments will compare methods on the same task: one ratio,
one capacity, or the full positive path. The worst-case bounds of parametric
flow cannot be used as a bound for the number of \FS{} pivots.

\bibliographystyle{plainnat}
\bibliography{references}
\end{document}
